% !TEX root = ../matan.tex

\section{Дифференцируемость высших порядков. Формула Тейлора}

\subsection*{Дифференцируемость высших порядков}

\begin{Definition}{$n$-кратно дифференцируемая функция}{}
	Пусть $f\colon\Omega\to\RR$, $\Omega\subset\RR^d$, $x$ --- внутренняя точка $\Omega$.
	\begin{itemize}
		\item $f$ называется \textbf{дважды дифференцируемой} в точке $x$, если все
		первые частные производные $\partial_j f$ ($j=1,\dots,d$) определены в
		окрестности $x$ и дифференцируемы в точке $x$.
		\item $f$ называется \textbf{$n$ раз дифференцируемой} в точке $x$, если все
		частные производные $(n-1)$-го порядка
		\[
		\partial_{k_1}\partial_{k_2}\cdots\partial_{k_{n-1}} f,
		\qquad k_1,\dots,k_{n-1}\in\{1,\dots,d\},
		\]
		определены в окрестности $x$ и дифференцируемы в точке $x$.
	\end{itemize}
\end{Definition}

\begin{Remark}{Симметрия высших производных}{}
	Если $f$ дважды дифференцируема в точке $x$, то по теореме Юнга
	$\partial_j\partial_k f(x)=\partial_k\partial_j f(x)$ для всех $j,k$. Более общо,
	для $n$ раз дифференцируемой функции значение
	$\partial_{k_1}\cdots\partial_{k_n} f(x)$ не меняется при любой перестановке
	$\sigma$ индексов:
	\[
	\partial_{k_1}\cdots\partial_{k_n} f(x)
	=\partial_{k_{\sigma(1)}}\cdots\partial_{k_{\sigma(n)}} f(x).
	\]
\end{Remark}

\subsection*{Дифференциалы высших порядков}

Пусть $f$ дважды дифференцируема в точке $x$. Дифференцируя разложение
$\partial_k f(\xi)=\partial_k f(x)+\sum_{j}\partial_j\partial_k f(x)(\xi_j-x_j)+o(|\xi-x|)$
и подставляя в первый дифференциал, приходим к квадратичной форме --- \textbf{второму
дифференциалу}. Обозначая приращение $dx_k=\xi_k-x_k$,
\[
d^2 f(x)=\sum_{k=1}^d\sum_{j=1}^d\frac{\partial^2 f}{\partial x_k\,\partial x_j}(x)\,dx_k\,dx_j;
\]
это симметричное билинейное отображение $\RR^d\times\RR^d\to\RR$.

В общем случае для $n$ раз дифференцируемой функции \textbf{дифференциал $N$-го порядка}
есть $N$-линейное отображение
\[
d^N f(x)(\eta^1,\dots,\eta^N)
=\sum_{k_1=1}^d\cdots\sum_{k_N=1}^d
\frac{\partial^N f}{\partial x_{k_1}\cdots\partial x_{k_N}}(x)\,
\eta^1_{k_1}\cdots\eta^N_{k_N},
\qquad \eta^1,\dots,\eta^N\in\RR^d.
\]
В формуле Тейлора все аргументы совпадают, $\eta^i=\xi-x$, и используется
сокращение $d^N f(x)(\xi-x,\dots,\xi-x)$.

\subsection*{Многочлен Тейлора}

Для $n$ раз дифференцируемой в точке $x$ функции \textbf{многочленом Тейлора}
степени $N\le n$ называется
\[
P_{f,N}(\xi)=\sum_{m=0}^{N}\frac{1}{m!}\,d^m f(x)(\xi-x,\dots,\xi-x)
=\sum_{m=0}^{N}\frac{1}{m!}\sum_{k_1,\dots,k_m=1}^{d}
\partial_{k_1}\cdots\partial_{k_m} f(x)\,(\xi_{k_1}-x_{k_1})\cdots(\xi_{k_m}-x_{k_m}).
\]
Ключевое свойство: все частные производные $P_{f,N}$ до порядка $N$ включительно в
точке $x$ совпадают с соответствующими производными $f$.

\subsection*{Формула Тейлора с остаточным членом в форме Пеано}

\begin{Theorem}{Формула Тейлора (форма Пеано)}{}
	Пусть $f\colon\Omega\to\RR$ дифференцируема $N$ раз в точке $x\in\Omega$. Тогда
	\[
	f(\xi)=P_{f,N}(\xi)+o(|\xi-x|^N),\qquad \xi\to x.
	\]
\end{Theorem}

\begin{proof}
	Индукция по $N$.

	\textbf{База} ($N=1$). Утверждение --- это определение дифференцируемости:
	$f(\xi)=f(x)+d_xf(\xi-x)+o(|\xi-x|)$.

	\textbf{Шаг} ($N-1\Rightarrow N$). Положим $r(\xi):=f(\xi)-P_{f,N}(\xi)$. Так как
	производные $P_{f,N}$ до порядка $N$ в точке $x$ совпадают с производными $f$,
	\[
	r(x)=0,\qquad
	(\partial_{j_1}\cdots\partial_{j_m} r)(x)=0\quad\text{для всех } m\le N.
	\]
	В частности, каждая первая производная $\partial_j r$ дифференцируема $(N-1)$ раз в
	точке $x$, и все её производные до порядка $(N-1)$ в точке $x$ равны нулю, поэтому
	её многочлен Тейлора степени $(N-1)$ тождественно нулевой. По индукционному
	предположению, применённому к $\partial_j r$,
	\[
	(\partial_j r)(\eta)=o(|\eta-x|^{N-1}),\qquad \eta\to x.
	\tag{$\ast$}
	\]

	Запишем приращение $r(\xi)-r(x)=r(\xi)$ телескопической суммой, изменяя координаты
	от $x$ к $\xi$ по одной:
	\[
	r(\xi)=\sum_{j=1}^{d}\Bigl[
	r(x_1,\dots,x_{j-1},\xi_j,\xi_{j+1},\dots,\xi_d)
	-r(x_1,\dots,x_{j-1},x_j,\xi_{j+1},\dots,\xi_d)\Bigr].
	\]
	К каждой скобке (отличающейся лишь в $j$-й координате) применим теорему Лагранжа по
	этой переменной: существует $\theta_j$ между $x_j$ и $\xi_j$ такое, что скобка равна
	$\partial_j r(\eta^{(j)})\cdot(\xi_j-x_j)$, где
	$\eta^{(j)}=(x_1,\dots,x_{j-1},\theta_j,\xi_{j+1},\dots,\xi_d)$. Точка $\eta^{(j)}$
	лежит между $x$ и $\xi$, поэтому $|\eta^{(j)}-x|\le|\xi-x|$. Тогда
	\[
	r(\xi)=\sum_{j=1}^{d}\partial_j r(\eta^{(j)})\,(\xi_j-x_j).
	\]
	Используя $(\ast)$ и $|\xi_j-x_j|\le|\xi-x|$, оценим каждое слагаемое:
	\[
	\bigl|\partial_j r(\eta^{(j)})\,(\xi_j-x_j)\bigr|
	\le o\bigl(|\eta^{(j)}-x|^{N-1}\bigr)\cdot|\xi-x|
	=o\bigl(|\xi-x|^{N-1}\bigr)\cdot \bigO(|\xi-x|)
	=o\bigl(|\xi-x|^{N}\bigr).
	\]
	Суммируя по $j=1,\dots,d$ (конечное число слагаемых), получаем
	$r(\xi)=o(|\xi-x|^N)$, что и требовалось.
\end{proof}

\subsection*{Формула Тейлора с остаточным членом в форме Лагранжа}

\begin{Theorem}{Формула Тейлора (форма Лагранжа)}{}
	Пусть $f\colon\Omega\to\RR$ дифференцируема $(N+1)$ раз в некоторой выпуклой
	окрестности $U_x\subset\Omega$ точки $x$. Тогда для любой точки $\xi\in U_x$
	найдётся точка $\theta$ на отрезке между $x$ и $\xi$, такая что
	\[
	f(\xi)=P_{f,N}(\xi)+\frac{1}{(N+1)!}\,d^{N+1}f(\theta)(\xi-x,\dots,\xi-x).
	\]
\end{Theorem}

\begin{proof}
	Сведём задачу к одномерной. Так как $U_x$ открыта и выпукла, отрезок
	$\{x+t(\xi-x):t\in[0,1]\}$ содержится в $U_x$, и функцию
	\[
	g(t):=f\bigl(x+t(\xi-x)\bigr)
	\]
	можно рассматривать на интервале $(-\eps,1+\eps)$. Как композиция $(N+1)$ раз
	дифференцируемой $f$ и аффинного отображения $t\mapsto x+t(\xi-x)$, функция $g$
	$(N+1)$ раз дифференцируема. По правилу цепочки последовательно получаем
	\[
	g^{(k)}(t)=\sum_{k_1=1}^d\cdots\sum_{k_k=1}^d
	\partial_{k_1}\cdots\partial_{k_k} f\bigl(x+t(\xi-x)\bigr)
	(\xi_{k_1}-x_{k_1})\cdots(\xi_{k_k}-x_{k_k})
	=d^k f\bigl(x+t(\xi-x)\bigr)(\xi-x,\dots,\xi-x).
	\]
	Применим к $g$ одномерную формулу Тейлора (Маклорена) с остаточным членом в форме
	Лагранжа на отрезке $[0,1]$: существует $s\in(0,1)$ такое, что
	\[
	g(1)=\sum_{k=0}^{N}\frac{g^{(k)}(0)}{k!}+\frac{g^{(N+1)}(s)}{(N+1)!}.
	\]
	Поскольку $g(1)=f(\xi)$, $g(0)=f(x)$ и $g^{(k)}(0)=d^k f(x)(\xi-x,\dots,\xi-x)$,
	первые $N+1$ слагаемых дают в точности $P_{f,N}(\xi)$. Полагая
	$\theta=x+s(\xi-x)$ (точка между $x$ и $\xi$), получаем
	\[
	f(\xi)=P_{f,N}(\xi)+\frac{1}{(N+1)!}\,d^{N+1}f(\theta)(\xi-x,\dots,\xi-x).\qedhere
	\]
\end{proof}
